% Objective puzzle check, not an empirical account of the solver.
\section{A Geometric Check for the Nine-Dot Example}
\label{sec:p19-insight-annex}

\begin{annexitem}{annex:p19.insight-nine-dot}{thm:p19-insight-nine-dot-path}{Four connected segments through the grid}
\begin{proof}
The first segment runs from \((0,0)\) to \((2,2)\) on \(y=x\);
it contains \((0,0),(1,1),(2,2)\).
The second runs from \((2,2)\) to \((2,-1)\) on \(x=2\);
it contains the remaining right-column points \((2,1),(2,0)\).
The third runs from \((2,-1)\) to \((-1,2)\) on \(y=1-x\);
it contains \((1,0),(0,1)\).
The fourth runs from \((-1,2)\) to \((2,2)\) on \(y=2\);
it contains \((0,2),(1,2)\) and the already visited \((2,2)\).
This list includes every point of \(\{0,1,2\}\times\{0,1,2\}\).

Adjacent segments have the same endpoint/startpoint by their displayed
order, so they form a continuous four-segment path without a pencil lift.
Their supporting lines are distinct, with diagonal-positive, vertical,
diagonal-negative and horizontal directions. Distinct supporting lines
can share at most a point, not a positive-length interval; no segment
is therefore retraced. Repeated passage through the endpoint \((2,2)\)
does not retrace a segment and is not forbidden by the stated instruction.
The two outside turning points are unrestricted by that instruction.
This proves only correctness of the specified geometric candidate.
\end{proof}
\end{annexitem}
